\documentclass{stex}

\libusepackage{naproche}
\libusepackage{copyright}
\libinput{preamble}

\begin{document}
\begin{smodule}{forthel}

\ifinputref\else\setsectionlevel{section}\fi

\begin{sfragment}[id=sec:forthel]{\ForTheL}
  \textsymdecl{forthel}{\ForTheL}
  \textsymdecl{asciidialect}{ASCII dialect}
  \textsymdecl{texdialect}{\TeX{} dialect}
  \textsymdecl{stexdialect}{\sTeX{} dialect}
  %
  \begin{sparagraph}[style=symdoc,for={forthel,asciidialect,texdialect,stexdialect}]
    The input language of \Naproche is called \defnotation\forthel
    (``\definiendum{forthel}{\textbf{For}mula \textbf{The}ory \textbf{L}anguage}'').
    \forthel comes in three dialects:
    \begin{enumerate}
      \item The \defnotation\asciidialect (which \Naproche expects when it reads a
        file whose file name ends with ``\texttt{.ftl}'') is a simple plain text format.
        Its name comes from the restriction that only characters from the Basic Latin
        Unicode block -- i.e.\ US-ASCII characters -- are allowed in \texttt{.ftl} files.
      \item The \defnotation\texdialect (which \Naproche expects when it reads a
        file whose file name ends with ``\texttt{.ftl.tex}'') is a \LaTeX-compatible
        format. This means that any document written in the \texdialect can directly
        be passed to a \TeX{} engine (e.g.\ \texttt{pdflatex} or \texttt{rustex}) to be
        converted to, e.g., PDF or HTML.
      \item The \defnotation\stexdialect (which \Naproche expects when it reads a
        file whose file name ends with ``\texttt{.ftl.en.tex}'') is an
        \textbf{experimental} \sTeX-compatible format which produces semantically
        annotated PDF or HTML documents when processed via \texttt{pdflatex} or
        \texttt{rustex}.
        The long-term goal is to automatically convert \forthel texts written in the
        \sr{asciidialect}{ASCII} or \texdialect to the \stexdialect so that
        texts in this format neither need to be written manually nor parsed by \Naproche.
    \end{enumerate}
  \end{sparagraph}
\end{sfragment}

\ifinputref\else\printlicense[CcByNcSa]{Marcel Schütz (2026)}\fi

\end{smodule}
\end{document}
